%% ma: L10-B13-0088
%% lop: 10
%% chuong: 5
%% bai: 13
%% dang: 10.13.15
%% loai: TN4
%% muc: TH
%% dap_an: "A"
%% nguon: "chuyen-de-cac-so-dac-trung-cua-mau-so-lieu-khong-ghep-nhom-toan-10-chuong-trinh-moi.pdf (GV. Phan Nhật Linh) (P03), Bài 02 Dạng 2 Phần I Câu 14"
%% kien_thuc: "Lớp 10 Bài 13 (số trung bình, trung vị, ý nghĩa)"
%% trang_thai: chinh-thuc
%% ghi_chu: "Lời giải gốc lập luận theo cỡ mẫu nhỏ; kho lập luận theo giá trị bất thường. Tự duyệt theo yêu cầu giáo viên 22/10/2026."
\begin{ex}
Tiền lương hằng tháng của $7$ nhân viên trong một công ty du lịch lần lượt là: $6{,}5$; $8{,}4$; $6{,}9$; $7{,}2$; $2{,}5$; $6{,}7$; $3{,}0$ (đơn vị: triệu đồng). Số đặc trưng phù hợp nhất để đại diện cho tiền lương hằng tháng của $7$ nhân viên có giá trị là
\choice{\True $6{,}7$ triệu đồng}{$7$ triệu đồng}{$5{,}9$ triệu đồng}{$6$ triệu đồng}
\loigiai{Sắp xếp: $2{,}5$; $3{,}0$; $6{,}5$; $6{,}7$; $6{,}9$; $7{,}2$; $8{,}4$. Mẫu có hai giá trị nhỏ bất thường ($2{,}5$; $3{,}0$) kéo số trung bình xuống ($\approx5{,}9$), nên dùng trung vị $M_e=6{,}7$ triệu đồng. Chọn A.}
\end{ex}
